\documentclass[11pt,a4paper]{article}
\usepackage[margin=2.2cm]{geometry}
\usepackage{amsmath,amssymb}
\usepackage{booktabs,array,xcolor}
\usepackage{graphicx}
\usepackage{parskip}
\usepackage{enumitem}
\usepackage{hyperref}
\hypersetup{colorlinks=true,linkcolor=blue,citecolor=blue,urlcolor=blue}
\usepackage[margin=2.2cm]{geometry}
\usepackage{amsmath,amssymb}
\usepackage{algorithm}
\usepackage{algpseudocode}
\usepackage{hyperref}
\hypersetup{colorlinks=true,linkcolor=blue,citecolor=blue,urlcolor=blue}



\title{ESN--A2 Stochastic-Volatility Model:\\
Architecture and hyperparameters tuning}
\author{Othmane Zarhali}
\date{July 14, 2026}

\begin{document}
\maketitle
\vspace{-0.3cm}
\noindent\rule{\textwidth}{0.8pt}

\begin{abstract}

This technical note is organised as follows:\\

\begin{enumerate}
    \item \textbf{Model specification.} introduces the complete ESN--A2 stochastic-volatility model, including its architecture, the data-generating processes, and the score function used for calibration.

    \item \textbf{Calibration framework.} presents the implementation choices required for the calibration procedure.

    \item \textbf{Scoring methodology.} defines the smooth Gaussian-kernel score that is employed throughout the study.

    \item \textbf{Model hyperparameters tuning and empirical results.} Section~\S4 evaluates the proposed architecture through several complementary analyses:
    \begin{enumerate}
        \item a complete grid search over the reservoir dimensions $(N_r,N_z)$, averaged over three independent random seeds (Section~\ref{sec:gridmatrix});
        \item the calibrated read-out parameters obtained for each volatility scenario (Section~\ref{sec:amplitude});
        \item the model's ability to reproduce the target Hurst exponent across four rough- and persistent-volatility scenarios (Section~\ref{sec:hurst}).
    \end{enumerate}

    \item \textbf{Conclusions.} Finally, it summarises the resulting optimal architecture and outlines the remaining directions for further validation.
\end{enumerate}
\end{abstract}

\tableofcontents
\vspace{0.4cm}

\section{The Model}

\subsection{Two parameter classes}

The pipeline separates two conceptually distinct parameter sets.

\begin{itemize}[itemsep=2pt]
\item \textbf{Hyperparameters} (shared across every DGP scenario,
found by an outer CMA-ES search): reservoir sizes $N_r,N_z$; z-bank
amplitude (\texttt{z\_strength}); pre-activation coefficients
$\beta_{\rm even},\beta_{\rm lin}$; damping $\gamma$
(\texttt{gamma\_norm}); local $z$-coupling $\delta_{\rm loc}$;
self-coupling scale $\nu_z$ (\texttt{zz\_scale}); sign probability
$\Pr(s_j=-1)$; orientation $\rho_r\in\{-1,+1\}$ (fixed to $-1$ by an
admissibility constraint, \S1.2.3). \textbf{7 hyperparameters in
total.}
\item \textbf{Calibrated parameters} (fit per DGP scenario inside an
inner Nelder--Mead loop): the reservoir's own memory exponent
$H_{\rm ESN}$ (distinct from the DGP's Hurst parameter, Remark 2), the
geometric OU-rate bounds $[\lambda_{\rm lo},\lambda_{\rm hi}]$, the
$z$-bank rate bounds $[a_{\rm lo},a_{\rm hi}]$, a bias offset
$\Delta b_0$, an output scale, the roughness-factor amplitude
$\rho_{\rm scale}$ (\texttt{rough\_scale}), a $z$-bank read-out profile
$(\texttt{zr\_lo},\texttt{zr\_hi})$, and a symmetric quadratic-term
matrix's two coefficients $(m_1,m_2)$. \textbf{12 calibrated
parameters in total.}
\end{itemize}

\subsection{ESN architecture}

The architecture has two coupled internal layers: a linear reservoir
$r$ (the ``rough factor'' bank) and a nonlinear, slower feedback bank
$z$ (the ``multiscale/Zumbach'' bank).

\subsubsection{The $r$-layer}

$r=(r_1,\ldots,r_{N_r})\in\mathbb{R}^{N_r}$ is a bank of $N_r$
discrete-time Ornstein--Uhlenbeck (OU) processes, all driven by the
same scalar innovation. Their mean-reversion rates are fixed,
geometrically spaced over the calibrated range
$[\lambda_{\rm lo},\lambda_{\rm hi}]$:
\begin{equation}
  \lambda_i = \lambda_{\rm lo}\Bigl(\frac{\lambda_{\rm hi}}{\lambda_{\rm lo}}\Bigr)^{\frac{i-1}{N_r-1}},
  \quad i=1,\ldots,N_r,
  \label{eq:geomspace}
\end{equation}
i.e.\ $\lambda_1=\lambda_{\rm lo}$ (slowest mode), $\lambda_{N_r}=\lambda_{\rm
hi}$ (fastest mode), uniformly spaced in $\log\lambda$ in between.
Given a single scalar innovation $\varepsilon_t\sim\mathcal N(0,1)$
drawn once per time step (the same draw feeds every component $i$), the
layer updates as
\begin{equation}
  r_{i,t} = e^{-\lambda_i\Delta t}\,r_{i,t-1} + \sqrt{1-e^{-2\lambda_i\Delta t}}\,\varepsilon_t,
  \quad i=1,\ldots,N_r,
  \label{eq:rlayer}
\end{equation}
with $\Delta t=1$ (one trading day). Sharing one common Brownian shock
across all $N_r$ modes turns the bank into a (Markovian,
multi-exponential) representation of a single rough process: any
linear combination $q^\top r_t$ is a finite sum-of-exponentials
approximation to a Volterra/fractional kernel driven by that same
shock, and the full vector $r_t$ also feeds a quadratic energy term
(\S1.2.3). Because each component is a stationary AR(1)/OU recursion,
its stationary variance is $\mathrm{Var}(r_i)=1$ for every $i$, and the
stationary cross-covariance between two components is
\begin{equation}
  \mathrm{Cov}(r_i,r_j) = \frac{b_ib_j}{\lambda_i+\lambda_j},\qquad b_i=\sqrt{2\lambda_i},
  \label{eq:crosscov}
\end{equation}
used to normalise the read-out vector, described next.

\paragraph{Read-out $q$.} The rough factor fed to the output is the
scalar projection $q^\top r_t$, where $q\in\mathbb{R}^{N_r}$ is built
from the calibrated exponent $H_{\rm ESN}$ as
\begin{equation}
  \tilde q_i = \lambda_i^{1/2-H_{\rm ESN}},\qquad
  q = \frac{\tilde q/\|\tilde q\|_2}{\sqrt{(\tilde q/\|\tilde q\|_2)^\top C(\tilde q/\|\tilde q\|_2)}},
  \label{eq:readout_q}
\end{equation}
so that $q^\top Cq=1$: the rough factor $q^\top r_t$ has unit
stationary variance for any choice of $H_{\rm ESN}$ or
$[\lambda_{\rm lo},\lambda_{\rm hi}]$.

\subsubsection{The $z$-layer}

$z=(z_1,\ldots,z_{N_z})\in\mathbb{R}^{N_z}$ is a second, separate bank
with its own geometrically spaced rates $a_j\in[a_{\rm lo},a_{\rm hi}]$,
\begin{equation}
  a_j = a_{\rm lo}\Bigl(\frac{a_{\rm hi}}{a_{\rm lo}}\Bigr)^{\frac{j-1}{N_z-1}},\quad j=1,\ldots,N_z,
  \label{eq:zrates}
\end{equation}
updated as a deterministic nonlinear filter of the $r$-layer's own
trajectory:
\begin{equation}
  z_{j,t} = e^{-a_j\Delta t}\,z_{j,t-1} + \bigl(1-e^{-a_j\Delta t}\bigr)\,m_t\,\zeta\,\tanh(u_{j,t}),
  \qquad
  m_t = \max\Bigl(0,\ 1-\gamma\,\|r_t\|_2/\sqrt{N_r}\Bigr),
  \label{eq:zupdate}
\end{equation}
with $\zeta=\texttt{z\_strength}$, $\gamma=\texttt{gamma\_norm}$, and
\begin{equation}
  u_{j,t} = \underbrace{s_jp_{1,j}p_{2,j}}_{\text{cross (Zumbach)}}
          + \underbrace{\beta_{\rm even}(0.7p_{1,j}^2+0.3p_{2,j}^2-1)}_{\text{even (variance-like)}}
          - \underbrace{\beta_{\rm lin}(0.7p_{1,j}+0.3p_{2,j})}_{\text{linear (leverage-like)}}
          + \underbrace{\nu_jz_{j,t-1}}_{\text{self}}
          + \underbrace{\tfrac12\delta_{\rm loc}(z_{j-1,t-1}+z_{j+1,t-1})}_{\text{local (chain) coupling}},
  \label{eq:preact}
\end{equation}
with boundary indices reflected, a fixed per-mode sign
$s_j\in\{-1,+1\}$ ($\Pr(s_j{=}+1)=\texttt{sign\_prob\_neg}$), and a
fixed per-mode self-coupling coefficient
$\nu_j\sim\mathrm{Unif}(-\nu_z,\nu_z)$ ($\nu_z=\texttt{zz\_scale}$).
The two reservoir taps $p_{1,j}=r_{f_i(j)}$, $p_{2,j}=r_{s_i(j)}$ move
in opposite directions across the spectrum as $j$ increases:
\begin{align}
  f_i(j) &= \min\Bigl(N_r{-}1,\ \tfrac{N_r}{2}+\bigl\lfloor(\tfrac{N_r}{2}{-}1)\tfrac{j}{N_z-1}\bigr\rfloor\Bigr)
  \quad(\tfrac{N_r}{2}\to N_r{-}1: \text{increasingly fast tap}), \label{eq:fi}\\
  s_i(j) &= \bigl\lfloor(\tfrac{N_r}{2}{-}1)\tfrac{N_z-1-j}{N_z-1}\bigr\rfloor
  \quad(\tfrac{N_r}{2}{-}1\to0: \text{increasingly slow tap}), \label{eq:si}
\end{align}
for $j=0,\ldots,N_z{-}1$ (0-indexed here to match the implementation).
Thus $z_0$ combines two adjacent, mid-spectrum $r$-modes, while
$z_{N_z-1}$ combines the single fastest mode with the single slowest
mode: the $z$-bank is explicitly constructed to read cross-timescale
(fast$\times$slow) products and sums out of the $r$-layer, and to fold
that information back into the volatility process through its own,
independently slow relaxation rates $a_j$ and the local chain coupling
$\delta_{\rm loc}$.

\subsubsection{Coupling to the output}
\label{sec:coupling}

The instantaneous log-volatility pre-activation reads both layers
additively, through three terms:
\begin{equation}
  \eta_t = b_0 + \rho_r\,\texttt{rough\_scale}\,(q^\top r_t)
         + \sum_{j=1}^{N_z} w_{j,z}\,z_{j,t}
         + r_t^\top Q\,r_t,
  \qquad
  \sigma_t = \sqrt{\mathrm{SIG\_MIN}^2+\mathrm{softplus}(\eta_t)^2},
  \label{eq:coupling}
\end{equation}
with
\begin{equation}
\begin{cases}
    w_{j,z} = \frac{\texttt{z\_readout}_j}{\sqrt{N_z}}\\
    \texttt{z\_readout}_j = \mathrm{geomspace}(\texttt{zr\_lo},\texttt{zr\_hi},N_z)[j]\\
     Q = \frac{1}{N_r}\,M,\qquad M = m_1 I_{N_r} + m_2\,qq^\top.
\end{cases}
  \label{eq:profile_and_Q}
\end{equation}
Daily log-returns are $dX_t=(\mu-\tfrac12\sigma_t^2)+\sigma_t\varepsilon_t$
(reusing the same $\varepsilon_t$ that drove $r_t$, Eq.~\eqref{eq:rlayer}),
so the return and the volatility innovation are perfectly
instantaneously correlated -- the source of the model's leverage
effect. Admissibility requires
$\kappa_0=\rho_r\,\texttt{rough\_scale}\,q^\top\sqrt{2\lambda}<0$
(negative leverage), enforced by fixing $\rho_r=-1$; since $q$ and
$\sqrt{2\lambda}$ are both strictly positive vectors, $\kappa_0<0$
holds automatically whenever $\texttt{rough\_scale}>0$. This
constraint depends only on the linear rough term, not on the $z$-bank
profile or the quadratic term.

\paragraph{The three read-out terms, in words.} (i) A linear rough
factor $\texttt{rough\_scale}\cdot q^\top r_t$: the r-layer's own,
Volterra-kernel-shaped contribution. (ii) A per-mode $z$-bank sum
$\sum_j w_{j,z}z_{j,t}$, with weights $\texttt{z\_readout}_j$ following
a smooth geometric \emph{profile} across the $N_z$ modes (2 calibrated
parameters $\texttt{zr\_lo},\texttt{zr\_hi}$, rather than one shared
scalar or $N_z$ fully free weights) -- since mode $j$ already carries a
fast$\leftrightarrow$slow meaning via the $f_i(j)/s_i(j)$ taps
(Eq.~\ref{eq:fi}--\ref{eq:si}), this profile lets different $z$-modes
matter differently to the volatility level. (iii) A quadratic-in-$r$
term $r_t^\top Qr_t=(1/N_r)(m_1\|r_t\|^2+m_2(q^\top r_t)^2)$, with
$Q=(1/N_r)M$ a genuine, \textbf{symmetric} matrix (not restricted to a
multiple of the identity): $m_1$ scales the $r$-layer's isotropic
instantaneous energy (already computed for the $z$-bank's gate $m_t$,
Eq.~\ref{eq:zupdate}), while $m_2$ scales a second, independent channel
that responds specifically to the \emph{squared} rough factor
$(q^\top r_t)^2$ -- the same direction that already drives the linear
rough term, now also entering quadratically.

\paragraph{Why this parameterisation (tractability).} A fully free
per-mode $z$-weight vector together with a fully free symmetric
$N_r\times N_r$ matrix $Q$ would add $N_z+N_r(N_r+1)/2$ calibrated
dimensions ($8+2080=2088$ at $N_r=64,N_z=8$; more at larger $N_r$) --
far beyond what a per-scenario Nelder--Mead multi-start can reliably
navigate, and likely to overfit the 11-term score against a handful of
DGP scenarios. $M=m_1I_{N_r}+m_2qq^\top$ is symmetric by construction
for \emph{any} real $(m_1,m_2)$ (a real linear combination of symmetric
matrices is symmetric); no positive-definiteness constraint is
imposed. $M$'s eigenvalues are $m_1$ (multiplicity $N_r-1$, on $q$'s
orthogonal complement) and $m_1+m_2\|q\|^2$ (multiplicity 1, along
$q$) -- exactly diagonal in the eigenbasis $\{q,q^\perp\}$, and either
eigenvalue may be positive, negative, or zero.

\paragraph{Remark 1 (Division of labour between the three read-out
terms).} The $r$-layer alone, through $q$ and $H_{\rm ESN}$, already
generates a rough-volatility-like process (a finite-dimensional
Markovian approximation of a fractional kernel). The $z$-layer adds
genuine nonlinearity (the $\tanh$ squashing and the quadratic/even
terms of Eq.~\eqref{eq:preact}) and an explicit, by-construction
fast$\times$slow cross-timescale interaction. The quadratic term
$r_t^\top Qr_t$ is a distinct, \emph{reservoir-wide} signal: unlike the
$z$-bank's own internal even/quadratic term, which only ever sees
\emph{two} taps at a time, $r_t^\top Qr_t$ reads the full
$N_r$-dimensional reservoir state, split by $M$'s construction into an
isotropic component ($m_1$) and a component aligned with the rough
read-out direction ($m_2$).

\paragraph{Remark 2 (Two different objects both called ``$H$'').}
The DGP's Hurst parameter $H_{\rm DGP}\in\{0.10,0.40\}$ (used to
generate the target paths) and the ESN's own memory-kernel exponent
$H_{\rm ESN}\in[0.02,0.49]$ (a calibrated hyperparameter of the
read-out weighting) are different objects. The inner calibration loop
searches over $H_{\rm ESN}$ so that the realised structure-function
estimate $\hat H_{\rm sim}$ of the simulated ESN path matches the
DGP's target $\hat H_c$; they are not required to coincide, and in
general do not (\S\ref{sec:hurst}).

\section{Data Generating Process (DGP) portfolio for the grid search procedure}

Four scenarios, equally weighted, define the target stylised facts:

\begin{center}
\begin{tabular}{@{}lll@{}}
\toprule
Name & Type & Parameters \\
\midrule
\texttt{fBM\_rough}   & \texttt{lnsv}    & $H=0.10,\ \lambda=1.0$ \\
\texttt{MRW\_taylor}  & \texttt{mrw}     & $\lambda^2=0.03,\ T_{\rm int}=252$ \\
\texttt{rBergomi\_Z}  & \texttt{bergomi} & $H=0.10,\ \eta=1.0,\ \rho=-0.70$ \\
\texttt{fBM\_persist} & \texttt{lnsv}    & $H=0.40,\ \lambda=1.5$ \\
\bottomrule
\end{tabular}
\end{center}

The \texttt{lnsv} (log-normal stochastic-volatility fBM) type sets
$\log\sigma_t=\log\sigma_0+\lambda B_t^H-\tfrac12\lambda^2t^{2H}$
(measuring $t$ in years, consistently, in both the fBM covariance and
the It\^o correction -- this is required for numerical stability at
$H=0.40$ over long simulated horizons, \S2); the \texttt{bergomi}
(rough Bergomi) type is the analogous
$V_t=\xi_0\exp(\eta B_t^H-\tfrac12\eta^2t^{2H})$ with correlated
leverage $\rho$; the \texttt{mrw} (multifractal random walk) type uses
a bounded, already-stationary log-covariance kernel
$\lambda^2\log(T_{\rm int}/|t-s|)$. For each scenario, $N_{\rm DGP\_PATHS}$
independent Monte-Carlo simulations of length $T_{\rm SIM}$ build a
reference vector of 11 stylised-fact centres/scales
(Table~\ref{tab:scoreterms}); the ESN, for a given architecture, is
calibrated to fit this reference.

\subsection{Score function}

Every stylised-fact term is expressed as a deviation
$d_k=w_k(1-f_k)\in[0,w_k]$ ($0=$ perfect, $w_k=$ worst), using a smooth
Gaussian-kernel proximity
\begin{equation}
  f(x,c,s) = \exp\Bigl(-\tfrac12\bigl(\tfrac{x-c}{s}\bigr)^2\Bigr),
  \label{eq:smoothkernel}
\end{equation}
rather than a piecewise-linear ``tent'' -- the smooth kernel provides a
never-zero gradient at any distance from target, whereas a tent kernel
$\max(0,1-|x-c|/s)$ has an exactly-zero gradient beyond one tolerance
width, which stalls a local optimiser (Nelder--Mead) once a statistic
drifts past tolerance (Figure~\ref{fig:kernels}). The total score is
\begin{equation}
  S = \sum_{k=1}^{11}w_k(1-f_k) + 5\cdot\mathrm{Stress}\ \in\ [0,15.3),
  \label{eq:score}
\end{equation}
where $\mathrm{Stress}\in\{0,1\}$ flags an inadmissible path (blow-up
or collapse of annualised volatility). $S=0$ requires a perfect match
on all 11 stylised facts and no stress event; $S$ is minimised directly
(no sign flip) by both the inner Nelder--Mead loop and the outer
CMA-ES search.

\begin{figure}[h]
\centering
\includegraphics[width=0.85\textwidth]{Numerical Protocols _ stylized facts reproduction/GridSearch_hyperparameters/kernel_shapes.png}
\caption{Tent vs.\ smooth (Gaussian) proximity kernel $f$ (left) and the
resulting per-term deviation $d=1-f$ (right). The tent is exactly flat
($d=1$, zero gradient) beyond one tolerance width $s$; the smooth
kernel never flattens.}
\label{fig:kernels}
\end{figure}

\begin{table}[h]
\centering
\caption{The 11 deviation terms $d_k=w_k(1-f_k)$ entering $S$ (Eq.~\eqref{eq:score}).}
\label{tab:scoreterms}
\begin{tabular}{@{}lp{7.2cm}r@{}}
\toprule
Statistic & Description & Weight $w_k$ \\
\midrule
$\hat H$ & Local H\"older exponent (structure function of $\log v$) & 2.2 \\
$\sigma_{\rm ann}$ & Mean annualised volatility & 1.4 \\
$\sigma^{0.995}_{\rm ann}$ & 99.5\% quantile of annualised vol & 0.6 \\
$\sigma^{\max}_{\rm ann}$ & Max annualised vol & 0.2 \\
$\rho_v$ & Mean variance autocorrelation (lags 1--20) & 0.8 \\
Taylor gap & $\rho_{|x|}-\rho_{x^2}$ & 1.0 \\
Taylor fraction & $\Pr(\rho_{|x|}>\rho_{x^2})$ & 0.8 \\
Zumbach & Past-return$\to$future-vol vs.\ past-vol$\to$future-return asymmetry & 1.0 \\
Leverage & $\mathrm{corr}(x_t,v_{t+L})$ & 1.1 \\
Kurtosis & Excess kurtosis of returns & 0.7 \\
$\max|\rho_x|$ & Max$|$return autocorrelation$|$ & 0.5 \\
\midrule
\multicolumn{2}{r}{$+\ 5\cdot\mathrm{Stress}$} & 5.0 \\
\midrule
\multicolumn{2}{r}{Total (worst case)} & 15.3 \\
\bottomrule
\end{tabular}
\end{table}

\subsection{Search procedure}

For a given architecture (the 7 hyperparameters of \S1.1 plus
$N_r,N_z$), the inner loop runs a multi-start Nelder--Mead per DGP
scenario to find
\[
H_{\rm ESN},\ \lambda_{\rm lo},\ \lambda_{\rm hi},\ a_{\rm lo},\ a_{\rm hi},\
\Delta b_0,\ \mathrm{scale},\ \texttt{rough\_scale},\
\texttt{zr\_lo},\ \texttt{zr\_hi},\ m_1,\ m_2
\]
(\textbf{12 parameters}) that minimise $S$ against that scenario's
reference; the outer loop averages $S$ across the 4 scenarios and
searches the 7 continuous hyperparameters via CMA-ES (Hansen \&
Ostermeier), for each $(N_r,N_z)$ cell of a discrete grid
(\S\ref{sec:gridmatrix}).

The inner loop uses \textbf{eleven} qualitatively different starting
points, kept whichever converges to the lowest $S$: a rough/fast
reservoir, a persistent/slow reservoir, an intermediate reservoir, a
wide-band/multiscale reservoir (spanning both day-scale and
multi-year OU modes simultaneously, so a single reservoir can support
both rough and persistent targets), a long-memory-extreme reservoir,
strong- and weak-amplitude starts (probing the $\texttt{rough\_scale}$/
$z$-profile range), and four quadratic-term starts (isotropic-positive,
isotropic-negative, $q$-aligned-positive, $q$-aligned-negative) probing
the two channels of $M$ independently.



\section{Grid-Matrix Hyperparameter Search: Algorithm}

The procedure described in \S1.3 (score function) and \S1.4 (search
procedure) is summarised below as three nested algorithms: the
top-level grid-matrix sweep with an outer CMA-ES search
(Algorithm~\ref{alg:gridsearch}), the per-scenario multi-start
Nelder--Mead inner loop it calls (Algorithm~\ref{alg:innercalibrate}),
and the smooth-kernel score evaluation both rely on
(Algorithm~\ref{alg:score}).

\begin{algorithm}[H]
\caption{ESN Grid-Matrix Hyperparameter Search}
\label{alg:gridsearch}
\begin{algorithmic}[1]
\Require DGP portfolio $\mathcal{D} = \{\texttt{fBM\_rough}, \texttt{MRW\_taylor}, \texttt{rBergomi\_Z}, \texttt{fBM\_persist}\}$
\Require Grid of reservoir sizes $\mathcal{G} = \{(N_r, N_z)\}$
\Require Multi-start set $\mathcal{P}_0$ of 11 starting points (Algorithm~\ref{alg:innercalibrate})
\Require CMA-ES budget (population size, number of generations)
\Ensure Optimal architecture $(N_r^\ast, N_z^\ast, \boldsymbol{\theta}^\ast)$ with minimal mean score $S^\ast$

\Statex
\Statex \textit{Stage 1 --- build the DGP reference (once, shared across the whole sweep)}
\For{each scenario $d \in \mathcal{D}$}
    \State Simulate $N_{\rm DGP\_PATHS}$ independent paths of length $T_{\rm SIM}$ under scenario $d$
    \State $\mathrm{ref}_d \gets$ 11 stylised-fact centres/scales (Table~\ref{tab:scoreterms})
\EndFor

\Statex
\Statex \textit{Stage 2 --- sweep the $(N_r,N_z)$ grid}
\State $S^{\rm best}_{\rm grid} \gets +\infty$
\For{each $(N_r, N_z) \in \mathcal{G}$}
    \State Initialise CMA-ES over the 7 shared hyperparameters $\boldsymbol{\theta} \in \mathbb{R}^7$
    \While{CMA-ES budget not exhausted}
        \For{each candidate $\boldsymbol{\theta}$ in the current CMA-ES generation}
            \For{each scenario $d \in \mathcal{D}$}
                \State $S_d \gets \Call{InnerCalibrate}{\boldsymbol{\theta}, N_r, N_z, \mathrm{ref}_d, \mathcal{P}_0}$
            \EndFor
            \State $\bar S(\boldsymbol{\theta}) \gets \dfrac{1}{|\mathcal{D}|}\sum_{d \in \mathcal{D}} S_d$ \Comment{mean score across the 4 scenarios}
        \EndFor
        \State Update CMA-ES using fitness values $\{\bar S(\boldsymbol{\theta})\}$ of the generation
    \EndWhile
    \State $(\boldsymbol{\theta}^\ast_{N_r,N_z},\, S^\ast_{N_r,N_z}) \gets$ best $(\boldsymbol{\theta}, \bar S(\boldsymbol{\theta}))$ found for this cell
    \If{$S^\ast_{N_r,N_z} < S^{\rm best}_{\rm grid}$}
        \State $S^{\rm best}_{\rm grid} \gets S^\ast_{N_r,N_z}$
        \State $(N_r^\ast, N_z^\ast, \boldsymbol{\theta}^\ast) \gets (N_r, N_z, \boldsymbol{\theta}^\ast_{N_r,N_z})$
    \EndIf
\EndFor
\State \Return $(N_r^\ast, N_z^\ast, \boldsymbol{\theta}^\ast, S^{\rm best}_{\rm grid})$
\end{algorithmic}
\end{algorithm}

\begin{algorithm}[H]
\caption{\textsc{InnerCalibrate} --- multi-start Nelder--Mead per-scenario calibration}
\label{alg:innercalibrate}
\begin{algorithmic}[1]
\Function{InnerCalibrate}{$\boldsymbol{\theta}, N_r, N_z, \mathrm{ref}, \mathcal{P}_0$}
    \State $S^{\rm best} \gets +\infty$, \ $\mathbf{p}^{\rm best} \gets \varnothing$
    \For{each starting point $\mathbf{p}_0 \in \mathcal{P}_0$}
        \Comment{11 qualitatively different starts: rough/fast, persistent/slow,
        intermediate, wide-band/multiscale, long-memory-extreme, strong-amplitude,
        weak-amplitude, isotropic-positive, isotropic-negative, $q$-aligned-positive,
        $q$-aligned-negative}
        \State $\mathbf{p}_0 = \bigl(H_0,\ \lambda_{{\rm lo},0},\ \lambda_{{\rm hi},0},\ a_{{\rm lo},0},\ a_{{\rm hi},0},\ \Delta b_{0,0},\ s_0,\ \rho_{{\rm scale},0},\ zr_{{\rm lo},0},\ zr_{{\rm hi},0},\ m_{1,0},\ m_{2,0}\bigr)$
        \State Let $\phi(\mathbf{p}) \gets \Call{Score}{\mathbf{p}, \boldsymbol{\theta}, N_r, N_z, \mathrm{ref}}$ \Comment{objective for this start}
        \State $\mathbf{p}^\ast \gets \Call{NelderMead}{\mathbf{p}_0,\ \phi}$
        \Comment{12-dimensional search: $H,\lambda_{\rm lo},\lambda_{\rm hi},a_{\rm lo},a_{\rm hi},\Delta b_0,\mathrm{scale},\rho_{\rm scale},zr_{\rm lo},zr_{\rm hi},m_1,m_2$}
        \If{$\phi(\mathbf{p}^\ast) < S^{\rm best}$}
            \State $S^{\rm best} \gets \phi(\mathbf{p}^\ast)$
            \State $\mathbf{p}^{\rm best} \gets \mathbf{p}^\ast$
        \EndIf
    \EndFor
    \State \Return $S^{\rm best}$ \Comment{return $\mathbf{p}^{\rm best}$ too if calibrated parameters are needed downstream}
\EndFunction
\end{algorithmic}
\end{algorithm}

\begin{algorithm}[H]
\caption{\textsc{Score} --- smooth-kernel data-adaptive score $S$}
\label{alg:score}
\begin{algorithmic}[1]
\Function{Score}{$\mathbf{p}, \boldsymbol{\theta}, N_r, N_z, \mathrm{ref}$}
    \State Build reservoir matrices from $\boldsymbol{\theta}$ (7 shared hyperparameters), $N_r, N_z$, and $\mathbf{p}$ (calibrated parameters)
    \If{admissibility constraint $\kappa_0 \geq 0$} \Comment{$\kappa_0 = \rho_r \cdot \rho_{\rm scale}\cdot q^\top\sqrt{2\lambda}$}
        \State \Return $S_{\max}$ \Comment{inadmissible path, worst-possible score}
    \EndIf
    \State Simulate the model path(s) under $\mathbf{p}, \boldsymbol{\theta}$
    \State Compute the 11 stylised facts $x_1,\dots,x_{11}$ from the simulated path(s) (Table~\ref{tab:scoreterms})
    \State $\mathrm{Stress} \gets \mathbb{1}[\text{blow-up or collapse of annualised volatility}]$
    \State $S \gets 0$
    \For{$k = 1,\dots,11$}
        \State $f_k \gets \exp\!\left(-\tfrac12\left(\dfrac{x_k - c_k}{s_k}\right)^{2}\right)$ \Comment{$c_k,s_k$ from $\mathrm{ref}$; Eq.~\eqref{eq:smoothkernel}}
        \State $S \gets S + w_k(1 - f_k)$ \Comment{$w_k$ from Table~\ref{tab:scoreterms}}
    \EndFor
    \State $S \gets S + 5\cdot \mathrm{Stress}$
    \State \Return $S$ \Comment{$S \in [0, 15.3)$, $S=0$ is a perfect match}
\EndFunction
\end{algorithmic}
\end{algorithm}


\section{Implementation Requirements}
\label{sec:bug1}

Two implementation details are required for the calibration to be
well-posed, and are stated here as requirements rather than as a
narrative of bugs found and fixed.

\paragraph{Time must be measured in years, consistently, in the fBM
DGP simulator.} The construction $\log\sigma_t=\log\sigma_0+\lambda
B_t^H-\tfrac12\lambda^2t^{2H}$ must use the same time unit in both the
fBM covariance and the It\^o correction term. Because $B_t^H\overset
{d}{=}t^HB_1^H$, the correction ($\propto t^{2H}$) dominates the
stochastic term ($\propto t^H$) for any fixed $\lambda>0$ as $t$ grows;
using a raw day-count index at $H=0.40$ over a 1200-day path drives
$\sigma_t$ to numerical zero (the correction reaches $t^{0.8}\approx291$
at $t=1200$, vs.\ a target of $O(1)$), and every downstream statistic
to NaN. Measuring time in years ($t_{\rm years}=t/252$) in both places
keeps $\sigma_t$ stationary and well-scaled around its target for
every $H$ in the portfolio.

\paragraph{The inner calibration loop requires multiple, qualitatively
different starting points.} The mapping from $H_{\rm ESN}$ to the
realised structure-function exponent is jointly determined by
$H_{\rm ESN}$, the OU-rate range $[\lambda_{\rm lo},\lambda_{\rm hi}]$,
and $N_r$ (Remark 2); a reservoir initialised with only fast modes
present injects high-frequency roughness that a single, locally
convergent Nelder--Mead run rarely escapes, regardless of which target
regime is being fit. The eleven-start scheme of \S1.5 -- in particular
the wide-band/multiscale start, whose modes span nearly six decades of
OU rate within one reservoir -- avoids this basin-of-attraction lock by
letting the score function decide, per scenario, how the read-out
$H_{\rm ESN}$ should weight each part of the spectrum, rather than the
start pre-committing to one extreme.

\section{Results}

\subsection{The optimal architecture}
\label{sec:optarch}

\begin{center}
\renewcommand{\arraystretch}{1.3}
\fbox{\begin{minipage}{0.9\textwidth}
\textbf{Optimal architecture:} $N_r=96$, $N_z=12$, with the 7 shared
hyperparameters of Table~\ref{tab:optarch}. \texttt{rough\_scale}, the
$z$-readout profile $(\texttt{zr\_lo},\texttt{zr\_hi})$, and the
quadratic-matrix coefficients $(m_1,m_2)$ are calibrated per DGP
scenario (\S\ref{sec:amplitude}), not fixed as part of this
architecture.
\end{minipage}}
\end{center}

\begin{table}[h]
\centering
\caption{The 7 shared hyperparameters of the optimal architecture.}
\label{tab:optarch}
\begin{tabular}{@{}lr@{}}
\toprule
Hyperparameter & Value \\
\midrule
$N_r$ & 96 \\
$N_z$ & 12 \\
\texttt{z\_strength} & 0.1929 \\
\texttt{even\_strength} & 2.2291 \\
\texttt{linear\_strength} & 0.1930 \\
\texttt{gamma\_norm} & 0.9651 \\
\texttt{local\_z\_strength} & 0.0476 \\
\texttt{zz\_scale} & 0.0397 \\
\texttt{sign\_prob\_neg} & 0.2219 \\
\texttt{rough\_orientation} & $-1.0$ (fixed) \\
\bottomrule
\end{tabular}
\end{table}

\textbf{Important caveat, stated plainly:} $N_r=96,N_z=12$ is the
argmin of the grid-matrix sweep below, but \S\ref{sec:gridmatrix}
shows this sweep cannot reliably distinguish between any of the 16
cells tested at the budget used here -- the cell-to-cell spread in
mean score is smaller than a single cell's own run-to-run noise. This
architecture should be treated as the current best point estimate, to
be confirmed or revised by a larger-budget re-run, not as a
definitively established optimum.

\subsection{Grid-matrix sweep}
\label{sec:gridmatrix}

Figure~\ref{fig:heatmap} evaluates the 7 shared hyperparameters of
Table~\ref{tab:optarch} at every cell of the production grid
$(N_r,N_z)\in\{32,48,64,96\}\times\{8,12,16,20\}$, holding them fixed
and only changing $N_r,N_z$ (each cell independently recalibrates all
12 inner-loop parameters per scenario), \textbf{averaged over 3
independent seeds per cell} (48 cell-evaluations total).

\begin{figure}[h]
\centering
\includegraphics[width=0.62\textwidth]{Numerical Protocols _ stylized facts reproduction/GridSearch_hyperparameters/esn_hyperparam_search_smooth_heatmap.png}
\caption{Score $S$ over the full $(N_r,N_z)$ grid, averaged over 3
seeds per cell.}
\label{fig:heatmap}
\end{figure}

\begin{table}[h]
\centering
\caption{Full $(N_r,N_z)$ score matrix: mean $\pm$ std over 3 seeds.}
\label{tab:gridmatrix}
\resizebox{\textwidth}{!}{%
\begin{tabular}{@{}lrrrr@{}}
\toprule
 & $N_z=8$ & $N_z=12$ & $N_z=16$ & $N_z=20$ \\
\midrule
$N_r=32$ & $6.05\pm1.63$ & $6.07\pm1.62$ & $6.05\pm1.62$ & $6.09\pm1.65$ \\
$N_r=48$ & $6.08\pm1.60$ & $6.08\pm1.61$ & $6.08\pm1.60$ & $6.08\pm1.60$ \\
$N_r=64$ & $6.09\pm1.59$ & $6.39\pm1.36$ & $6.38\pm1.34$ & $6.37\pm1.34$ \\
$N_r=96$ & $6.38\pm1.36$ & $\mathbf{5.79\pm0.56}$ & $6.45\pm1.47$ & $6.45\pm1.47$ \\
\bottomrule
\end{tabular}%
}
\end{table}

\paragraph{Statistical read of the sweep.} $N_r=96,N_z=12$ has both the
lowest 3-seed mean ($5.79$) and the smallest standard deviation
($0.56$) of any cell. However: (i) the spread between the best and
worst cell means across the entire grid is only $0.66$ (from $5.79$ to
$6.45$), while (ii) the typical cell's own seed-to-seed standard
deviation is $1.46$ -- more than double the grid's entire cell-to-cell
spread. The difference between $N_r=96,N_z=12$ and the next most
plausible candidate, $N_r=64,N_z=8$, is $0.30$ against a standard error
of $0.97$ ($0.30/0.97\approx0.31$, far short of the $\gtrsim2$ needed
for even weak significance at $n=3$ per cell). \textbf{At this budget,
no $(N_r,N_z)$ cell in this grid is reliably distinguishable from any
other.} $N_r=96,N_z=12$ is reported as the architecture throughout the
remainder of this note because it is the argmin, with this caveat
attached to every result that follows.

\subsection{Per-scenario calibrated read-out parameters}
\label{sec:amplitude}

Table~\ref{tab:perscenario} reports the calibrated
\texttt{rough\_scale}, $z$-readout profile
$(\texttt{zr\_lo},\texttt{zr\_hi})$, and quadratic-matrix coefficients
$(m_1,m_2)$ at $N_r=96,N_z=12$, one independent set per DGP scenario
(single demonstration run).

\begin{table}[h]
\centering
\caption{Per-scenario calibrated read-out parameters at $N_r=96,N_z=12$.}
\label{tab:perscenario}
\begin{tabular}{@{}lccccc@{}}
\toprule
Scenario & \texttt{rough\_scale} & \texttt{zr\_lo} & \texttt{zr\_hi} & $(m_1,\ m_2)$ & Score $S$ \\
\midrule
\texttt{fBM\_rough}    & 0.616 & 0.201 & 0.015 & $(-0.784,\ +0.0002)$ & 4.601 \\
\texttt{MRW\_taylor}   & 0.372 & 0.044 & 0.081 & $(+0.0003,\ +0.0001)$ & 4.816 \\
\texttt{rBergomi\_Z}   & 0.414 & 0.197 & 0.019 & $(-0.610,\ -0.00005)$ & 5.581 \\
\texttt{fBM\_persist}  & 0.365 & 0.035 & 0.098 & $(-0.0001,\ -0.00003)$ & 5.569 \\
\bottomrule
\end{tabular}
\end{table}

\paragraph{The quadratic term is weakly and inconsistently identified
across scenarios.} \texttt{fBM\_rough} and \texttt{rBergomi\_Z} (both
$H=0.10$ targets) calibrate a large, same-signed $m_1$ ($-0.784$ and
$-0.610$); \texttt{MRW\_taylor} and \texttt{fBM\_persist} sit at
$m_1,m_2\approx0$. Repeated testing at other architectures shows this
pattern is \textbf{not stable}: at a different reservoir size, a
different pair of scenarios showed a large value in a different
coefficient ($m_2$, not $m_1$), and at yet another size, a different
pair again showed instability with standard deviation exceeding the
mean across seeds. \textbf{The most defensible reading is that the
quadratic channel absorbs a weak, largely interchangeable degree of
freedom that a single demonstration run cannot reliably pin down to
one coefficient, one scenario pairing, or one sign} -- not a genuine,
reproducible per-scenario mechanism. Treating $m_1,m_2\approx0$ as the
default working hypothesis, with occasional large excursions expected
from calibration noise rather than signal, is the current best
practice until confirmed otherwise by a multi-seed, increased-budget
re-run at this specific architecture.

\subsection{Hurst-matching accuracy}
\label{sec:hurst}

Table~\ref{tab:hurst} reports the local H\"older exponent match
$\hat H_{\rm sim}$ vs.\ each scenario's target $\hat H_c$, at
$N_r=96,N_z=12$ (single demonstration run, reduced calibration
budget).

\begin{table}[h]
\centering
\caption{$\hat H$ match at $N_r=96,N_z=12$ (single demonstration run, reduced budget).}
\label{tab:hurst}
\begin{tabular}{@{}lccc@{}}
\toprule
Scenario & $\hat H_{\rm sim}$ & $\hat H_c$ (target) & $\Delta\hat H$ \\
\midrule
\texttt{fBM\_rough}   & 0.167 & 0.092 & 0.076 \\
\texttt{MRW\_taylor}  & 0.529 & 0.364 & 0.165 \\
\texttt{rBergomi\_Z}  & 0.153 & 0.097 & 0.056 \\
\texttt{fBM\_persist} & 0.531 & 0.350 & 0.182 \\
\bottomrule
\end{tabular}
\end{table}

Match quality at this budget is uneven across scenarios
($\Delta\hat H$ ranging $0.056$--$0.182$). A separate, larger-budget
calibration run (longer simulated horizon, more Monte-Carlo paths per
scenario, more Nelder--Mead iterations per start, averaged over
several independent seeds) has previously been shown, at a different
architecture, to close this gap substantially for the scenarios most
affected (down to $\Delta\hat H\leq0.015$ for the worst-affected
scenarios in that test); the same increased-budget procedure has not
yet been re-applied at this note's $N_r=96,N_z=12$ architecture, and
doing so is the most direct outstanding validation step (\S5).

\begin{figure}[h]
\centering
\includegraphics[width=0.95\textwidth]{Numerical Protocols _ stylized facts reproduction/GridSearch_hyperparameters/esn_hyperparam_search_smooth_overview.png}
\caption{Six-panel overview at $N_r=96,N_z=12$: per-cell score sweep,
score-vs-$N_r$ distribution, full heatmap, per-DGP score,
hyperparameter values vs.\ the A2-981003 reference architecture, and
achieved vs.\ target $\hat H$.}
\label{fig:overview}
\end{figure}

\begin{figure}[h]
\centering
\includegraphics[width=0.95\textwidth]{Numerical Protocols _ stylized facts reproduction/GridSearch_hyperparameters/esn_hyperparam_search_smooth_deviation.png}
\caption{Score trajectory per $(N_r,N_z)$ cell and the 11-term deviation
breakdown at $N_r=96,N_z=12$.}
\label{fig:deviation}
\end{figure}

\section{Limitations and Recommended Steps}

\begin{enumerate}[itemsep=2pt]
\item \textbf{Confirm the architecture at production budget.} The
grid-matrix sweep (\S\ref{sec:gridmatrix}) cannot currently distinguish
$N_r=96,N_z=12$ from most of the grid, including $N_r=64,N_z=8$. A
larger-budget, multi-seed re-run (more seeds than 3, and/or a longer
per-seed calibration budget) is needed before this architecture is
treated as settled.
\item \textbf{Re-verify Hurst-matching accuracy at production budget.}
\S\ref{sec:hurst}'s $\Delta\hat H$ values are from a single,
reduced-budget run; an increased-budget, multi-seed re-run at
$N_r=96,N_z=12$ specifically has not yet been performed.
\item \textbf{Do not over-interpret any single-run $(m_1,m_2)$
pattern.} \S\ref{sec:amplitude} shows this coefficient pair is weakly
identified and its apparent ``pattern'' changes with architecture and
seed; treat $m_1,m_2\approx0$ as the default and require multi-seed
confirmation before attributing a large calibrated value to any
scenario-specific mechanism.
\item \textbf{A still richer quadratic term, if warranted.} The
current $Q=(1/N_r)(m_1I_{N_r}+m_2qq^\top)$ is a 2-eigenvalue
simplification of a fully free symmetric matrix, chosen for
tractability. If a future multi-seed test shows a genuine,
reproducible need for direction-specific quadratic sensitivity beyond
the $q$ direction, a natural extension is a small basis of $q$-like
vectors at several fixed roughness exponents $H_k$ (each
unit-$C$-normalised as in Eq.~\eqref{eq:readout_q}), each with its own
calibrated coefficient -- still symmetric by construction, still
tractable (one new scalar per basis vector).
\item Track both the aggregate score $S$ and its 11-term breakdown in
every calibration run: the breakdown is the fastest diagnostic for any
future architecture or read-out change.
\end{enumerate}

\end{document}
